\section{The single-block identity and the local majorant induction}
\label{tp-section}

We retain the product of local fields $K$, compact open self-dual
subring $R$, character $\psi$, and orbital kernel $\mathcal K$ of
Section~\ref{sec:theta-orbital-identity}. Haar measure gives $R$
volume one. Write $[a]$ for the class of $a$ in $K/R$, and also
write $\psi$ for the induced character of that quotient. For an
integral invertible matrix $Z$ of size $n$, put
\[
 C_Z=R^n/ZR^n,\qquad e_Z=[e_n],\qquad
 \lambda_Z([v])=[(Z^{-1}v)_n].
\]
In size zero the module, marked vector and functional are zero.
Isomorphisms of marked modules preserve both the vector and the
functional; an unqualified isomorphism below preserves only the
$R$-module structure. All sums over finite modules are ordinary
finite sums. For a block configuration $h$, define
\[
 c(h)=|\det Z(h)|_K^{-1},\qquad N_0(h,h')=\mathcal K(h,h'),
\]
\[
 N_s(h,h';t,t')=|tt'|_K^{s(s-1)/2}
           \mathcal K(h[s,t],h'[s,t'])\quad(s\ge1).
\]
Here $h[s,t]$ means that the block $tJ_s$ is appended to $h$,
where $J_s$ reverses the coordinates. These are the source's
normalizations\cite{alpo-uniform-source}; in particular $N_1$ has no additional modulus
factor. We prove the exact single-block identity and then deduce
the local majorant theorem from the orbital germ proved in
Theorem~\ref{ga-actual-germ}.

\subsection{Duality and a bordered cokernel}

\begin{lemma}\label{tp-dual}
For $C=R^n/ZR^n$ as above, put $C^*=\operatorname{Hom}_R(C,K/R)$.
Then $\#C^*=\#C=|\det Z|_K^{-1}$, and
\[
 \sum_{\alpha\in C^*}\psi(\alpha(x))
       =\#C\,1_{x=0},\qquad
 \sum_{x\in C}\psi(\alpha(x))=\#C\,1_{\alpha=0}.
\]
\end{lemma}

\begin{proof}
An $R$-linear functional on $R^n$ with values in $K/R$ is
represented by a row $\xi\in K^n$ modulo $R^n$. It vanishes on
$ZR^n$ exactly when $\xi Z\in R^n$. Consequently
$C^*=R^nZ^{-1}/R^n$ as an additive group. Multiplication by $Z$
identifies this with the quotient of $R^n$ by $R^nZ$, whose
order is $|\det Z|_K^{-1}$ by the Haar-volume calculation in
Lemma~\ref{th-coker}. The same calculation gives the order of $C$.

If $x=[v]\ne0$, then $Z^{-1}v$ has a coordinate outside $R$.
Self-duality of $R$ gives an $r\in R$ on which its pairing by
$\psi$ is nontrivial. The corresponding row $r e_i^tZ^{-1}$
represents an $\alpha\in C^*$ with $\psi(\alpha(x))\ne1$.
Translation by that $\alpha$ makes the first sum zero. If
$\alpha\ne0$, choose $x$ with $\alpha(x)\ne0$ in $K/R$.
Self-duality gives $r\in R$ with $\psi(\alpha(rx))\ne1$;
translation by $rx$ makes the second sum zero. In the zero cases
every summand is one. This proof uses no principal-ideal property.
\end{proof}

For rows $\xi$, columns $\eta$, and $c\in K^\times$, set
\begin{equation}\label{tp-border}
 B(Z,\xi,\eta,c)=
 \begin{pmatrix}Z&Z\eta\\ \xi Z&\xi Z\eta+c\end{pmatrix}.
\end{equation}
Suppose $Z$ is integral and invertible. The conditions
$\xi Z\in R^n$ and $Z\eta\in R^n$ identify $\xi\bmod R^n$
with $\alpha\in C_Z^*$ and $\eta\bmod R^n$ with
$\beta=[Z\eta]\in C_Z$. The remaining condition for
\eqref{tp-border} to be integral is
\begin{equation}\label{tp-integral-condition}
                         \alpha(\beta)=-[c].
\end{equation}
Every pair $(\alpha,\beta)$ has such representatives. These
assertions follow either by applying $Z^{-1}$, or by evaluating
the functional on the standard basis of $R^n$.

\begin{lemma}\label{tp-border-isomorphisms}
Let $Z,Z'$ be integral invertible matrices, possibly of different
sizes. Fix $c\in K^\times$ and pairs $(\alpha,\beta)$ and
$(\alpha',\beta')$ satisfying \eqref{tp-integral-condition}.
There is a bijection between marked isomorphisms of the two
bordered cokernels and the unqualified isomorphisms
$\theta:C_Z\to C_{Z'}$ satisfying
\[
                 \alpha'\theta=\alpha,\qquad
                 \theta\beta=\beta'.
\]
\end{lemma}

\begin{proof}
For $C=C_Z$, form the $R$-module
\[
 E_\alpha=\{(x,k)\in C\times K:[k]=-\alpha(x)\}.
\]
The map $(v,r)\mapsto([v],r-\xi v)$ from $R^n\oplus R$
onto $E_\alpha$ has kernel consisting exactly of
$(Zu,\xi Zu)$ with $u\in R^n$. Surjectivity follows by lifting
$x$ to $v$ and using $[k]=-\alpha(x)$ to see that $k+\xi v\in R$.
Thus it identifies the quotient by the first $n$ columns of
\eqref{tp-border} with $E_\alpha$. The last column maps to
$g=(\beta,c)$, and the last standard vector maps to $e=(0,1)$.
Block inversion gives the last inverse coordinate
$c^{-1}(r-\xi v)$. Hence the bordered marked module is
\[
 E_\alpha/Rg,\qquad [e],\qquad [(x,k)]\longmapsto[k/c].
\]
Here $g\in E_\alpha$ by \eqref{tp-integral-condition}, and
$Rg$ is a free copy of $R$ because $c\in K^\times$.

An isomorphism of these marked quotients lifts uniquely to an
isomorphism $E_\alpha\to E_{\alpha'}$ preserving $k$. Indeed,
given a representative of the image class, preservation of the
functional says that its second coordinate differs from $k$ by
an element of $cR$. Subtracting the corresponding multiple of
$g'=(\beta',c)$ gives a unique lift with second coordinate $k$.
Uniqueness proves additivity and $R$-linearity, and the same
construction for the inverse proves bijectivity. The lift sends
$g$ to $g'$ and $e$ to $e'$, by their quotient classes and their
respective second coordinates $c$ and $1$.

Quotienting $E_\alpha$ by $Re$ recovers $C$, with functional
$-\alpha$ determined by $k\bmod R$. The lifted isomorphism
therefore induces $\theta$ with $\alpha'\theta=\alpha$ and
$\theta\beta=\beta'$. Conversely such a $\theta$ gives
$(x,k)\mapsto(\theta x,k)$ on the two $E$'s. It preserves
$e,g$ and descends to a marked isomorphism of the bordered
quotients. These constructions are inverse, proving the claim.
\end{proof}

\subsection{The exact single-block sum}

For marked modules $(C,e,\lambda)$ and $(C',e',\lambda')$,
define
\begin{equation}\label{tp-finite-inner}
 I_c(C,C')=
 \sum_{\theta:C\simeq C'}\ 
 \sum_{\substack{\alpha\in C^*,\ \beta\in C\\
                         \alpha(\beta)=-[c]}}
 \psi\bigl(\alpha(e-\theta^{-1}e')
                    +\lambda(\beta)-\lambda'(\theta\beta)\bigr).
\end{equation}
The outer isomorphisms are unqualified. The finite expression
depends only on $c\bmod R$ and is zero for all but finitely many
such classes.

\begin{proposition}\label{tp-explicit-single}
For configurations $h,h'$ and $c\in K^\times$,
\begin{equation}\label{tp-single-formula}
 N_1(h,h';c,c)=
 \sum_{q\in\mathcal I_h}\sum_{q'\in\mathcal I_{h'}}
       a_q\overline{a_{q'}}\,I_c(C_{Z_q},C_{Z_{q'}}),
\end{equation}
where each cokernel on the right retains its original marked
vector and last inverse-coordinate functional.
\end{proposition}

\begin{proof}
Fix lower and upper representatives $(x,y)$ of an orbital class
for $h$, so that $Z=xZ(h)y$. Representatives over this class
after appending the size-one block are
\[
 x_+=\begin{pmatrix}x&0\\ \xi x&1\end{pmatrix},\qquad
 y_+=\begin{pmatrix}y&y\eta\\0&1\end{pmatrix}.
\]
The rows $\xi$ and columns $\eta$ range independently over
$K^n/R^n$. To verify this parametrization, project a lower or
upper representative to its upper left block. Use an integral
change to put that block equal to the chosen representative.
The remaining integral changes add an arbitrary integral row
to $\xi$, or an arbitrary integral column to $\eta$; conversely
only such changes keep the upper left representative fixed.
This proves both existence and uniqueness of the indicated
coset coordinates.

Matrix multiplication gives
$x_+\operatorname{diag}(Z(h),c)y_+=B(Z,\xi,\eta,c)$.
Since the last column of $x$ and last row of $y$ are the last
standard vector and covector, the additional phase is
\[
 \psi(\xi_n+\eta_n)
       =\psi\bigl(\alpha(e_Z)+\lambda_Z(\beta)\bigr).
\]
For size zero both extra terms are zero. Integrality is precisely
\eqref{tp-integral-condition}; if the bordered matrix is integral,
its upper left matrix $Z$ is integral as well. Lemma~\ref{tp-border-isomorphisms}
now replaces the marked-isomorphism count by the number of
$\theta$ satisfying $\alpha'=\alpha\theta^{-1}$ and
$\beta'=\theta\beta$. The difference of the two phases is
exactly the argument of $\psi$ in \eqref{tp-finite-inner}.
Summing over the remaining $\alpha,\beta$ proves
\eqref{tp-single-formula}. The base orbital sets and all modules
are finite, so every rearrangement is finite.
\end{proof}

\begin{theorem}\label{tp-parseval}
The function $c\mapsto N_1(h,h';c,c)$ is $R$-periodic on
$K^\times$ and extends to a finitely supported function on
$K/R$. For any choice of unit representative $c_x$ of each
additive class $x\in K/R$,
\begin{equation}\label{tp-parseval-eq}
 \sum_{x\in K/R}N_1(h,h';c_x,c_x)
                       =c(h)c(h')N_0(h,h').
\end{equation}
\end{theorem}

\begin{proof}
Periodicity and finite support follow from
\eqref{tp-finite-inner} and \eqref{tp-single-formula}.
Each additive class has a representative in $K^\times$:
adding a suitable element of the open subgroup $R$ avoids
the finitely many coordinate hyperplanes of the product $K$.
For a fixed $\theta$ in \eqref{tp-finite-inner}, summing over
$c\bmod R$ removes the condition $\alpha(\beta)=-[c]$, since
each pair $(\alpha,\beta)$ determines exactly one such class.
The remaining sum factors as
\[
 \left(\sum_{\alpha\in C^*}
       \psi(\alpha(e-\theta^{-1}e'))\right)
 \left(\sum_{\beta\in C}
       \psi\bigl((\lambda-\lambda'\theta)(\beta)\bigr)\right).
\]
Lemma~\ref{tp-dual} makes this $(\#C)^2$ if $\theta$ preserves
both marks, and zero otherwise. An isomorphism $C\simeq C'$
implies $\#C=\#C'$. Therefore
\[
 \sum_{c\bmod R}I_c(C,C')
       =\#C\,\#C'\,\#\operatorname{Iso}((C,e,\lambda),(C',e',\lambda')).
\]
Every integral orbital matrix for $h$ has cokernel order $c(h)$,
and similarly for $h'$. Substitution in \eqref{tp-single-formula}
gives \eqref{tp-parseval-eq}.
\end{proof}

\begin{example}\label{tp-factor-example}
The factors $c(h)c(h')$ cannot be deleted. Take
$K=\mathbb Q_2$, $R=\mathbb Z_2$, and $h=h'=[1,2]$.
There is just one base orbital matrix, namely $(2)$; its
marked cokernel is $C=\mathbb Z/2\mathbb Z$, with $e=1$
and $\lambda(x)=[x/2]$. It has only the identity module
automorphism. Write $\alpha_r(x)=[rx/2]$, $r=0,1$.
The phase in \eqref{tp-finite-inner} is always one.
Three pairs $(r,\beta)$ have $r\beta=0$ and one has
$r\beta=1$. Thus $N_1$ takes values $3$ on the zero coset,
$1$ on $1/2+\mathbb Z_2$, and zero elsewhere. Its sum is
$4=c(h)^2N_0(h,h)$, whereas $N_0(h,h)=1$.
This is a counterexample to an incorrectly normalized identity,
not to \eqref{tp-parseval-eq} or to the submitted theorem.
\end{example}

\subsection{Support and the unrestricted diagonal}

\begin{lemma}\label{tp-minor-support}
For $s\ge1$, a nonzero $N_s(h,h';t,t')$ implies
\[
                 (\det Z(h))t\in R,
                 \qquad(\det Z(h'))t'\in R.
\]
If both $Z(h)$ and its inverse are integral, and likewise for
$h'$, then $N_0(h,h')=1$.
\end{lemma}

\begin{proof}
Let $r$ be the total size of $h$ and $n=r+s$. In an integral
orbital matrix $x\operatorname{diag}(Z(h),tJ_s)y$, take the
minor in the first $r+1$ rows and columns $1,\ldots,r,n$.
Multiplication on the left restricts to a lower unipotent
matrix on these rows and has determinant one. Before that
multiplication, the last selected row is zero in the first $r$
columns and has last entry $t$; the upper left $r$ by $r$
block is $Z(h)$ times an upper unipotent matrix. The minor is
therefore exactly $(\det Z(h))t$. Integrality puts it in $R$.
If the kernel is nonzero, at least one integral orbital matrix
occurs on each side, proving both assertions of support.

For the last assertion every reversing block of $Z(h)$ is
invertible over $R$. An integral matrix $xZ(h)y$ recovers
integral $x,y$ by successive block elimination. Indeed, the
first diagonal block is $t_1J_{s_1}$ times an upper unipotent
matrix, so the latter is integral and invertible over $R$.
The other blocks in its row and column recover the first
upper and lower off-diagonal blocks integrally. Their Schur
complement is integral and repeats the argument on the
remaining blocks. Thus only the integral orbital coset
occurs. Its phase is one and its cokernel is zero. The same
holds for $h'$, and the zero marked modules have one
isomorphism. This also covers empty configurations.
\end{proof}

For equal-size configurations $h,h'$ put
\[
 U_m(h,h';t,t')=|tt'|_K^{m(m+1)/2}
   \sum_{c\bmod R}\mathcal K(h[m,t][1,c],h'[m,t'][1,c]),
\]
using unit representatives. Theorem~\ref{tp-parseval} applies
to the configurations $h[m,t]$ and $h'[m,t']$. Since
$c(h[m,t])=c(h)|t|_K^{-m}$, it gives the exact identity
\begin{equation}\label{tp-unrestricted}
             U_m(h,h';t,t')=c(h)c(h')N_m(h,h';t,t').
\end{equation}

Fix a finite collection $\mathcal F$ of equal-size configurations.
Write $D_s=\operatorname{diag}_{h\in\mathcal F}(c(h)^s)$,
and regard kernels as matrices on $\mathcal F$.
Suppose a positive $R$-periodic majorant $\Theta_m$ satisfies
\[
 N_m(t,t)\le\Theta_m(t),\qquad
 \sum_{x\in K/R}\Theta_m(x)\le D_mN_0D_m,
\]
with finite support on $K/R$. Each positive summand is at most
the whole sum. Conjugating by $D_1$ and using
\eqref{tp-unrestricted} proves
\begin{equation}\label{tp-diagonal}
                     U_m(t,t)\le D_{m+1}N_0D_{m+1}
                         \quad(t\in K^\times).
\end{equation}
In particular the same bound holds for every convex average
of these unrestricted diagonals. No new analytic estimate is
required for this step.

\subsection{Completion of the induction from the germ}

\begin{theorem}\label{tp-majorant}
For every $s\ge1$ and finite family $\mathcal F$ of equal-size
configurations there is an $R$-periodic matrix-valued function
$\Theta_s$ with
\begin{equation}\label{tp-majorant-properties}
 0\le N_s(t,t)\le\Theta_s(t)\quad(t\in K^\times),\qquad
 \sum_{x\in K/R}\Theta_s(x)\le D_sN_0D_s.
\end{equation}
It is zero outside the fixed compact open $R$-module
\[
                 L_c=R+\sum_{h\in\mathcal F}(\det Z(h))^{-1}R.
\]
Its $(h,h')$ entry vanishes when
$|\det Z(h)|_K\ne|\det Z(h')|_K$.
If all $Z(h)$ are invertible over $R$, then
\[
                      0\le\Theta_s(x)\le1_R(x)N_0.
\]
The construction uses the germ uniformly on each fixed compact
parameter set. Theorem~\ref{ga-actual-germ} supplies exactly
that uniformity, independently of the comparison configuration.
\end{theorem}

\begin{proof}
For $s=1$, take $\Theta_1$ to be the periodic extension of
$N_1(t,t)$. Positivity is the Gram identity, its exact mass is
Theorem~\ref{tp-parseval}, and its support is
Lemma~\ref{tp-minor-support}. Determinant support gives the
claimed zero entries. This proves the base case.

Suppose the result is known for $s=m\ge1$. Fix a compact open
set $\Omega\subset K^\times$ invariant under $(m+1)$st roots
of unity. Choose additive lattices $L_1\subset R\subset L_2$
so that $\Omega,L_c\subset L_1^\perp$, and $\Omega$ and every
entry of $N_{m+1}(z,z)$ on it are invariant under translation
by $L_2^\perp$. Such lattices exist by compactness and the
uniform local constancy already proved in
Section~\ref{sec:theta-orbital-identity}; shrinking a lattice
to a power of $p$ times $R$ preserves these conditions.
Choose representatives $A$ of $L_2/R$, $B$ of $R/L_1$, and
put $a=\#A$, $b=\#B$.

For all shifts $\sigma=\rho+v$, with $\rho\in A$ and $v\in B$,
choose one sufficiently small $t\in K^\times$ and set
$t_\sigma^{-1}=t^{-1}+\sigma$. Their absolute values equal
$|t|_K$. The germ consequence in
Section~\ref{sec:theta-orbital-identity} identifies the truncated
Gram matrix by
\begin{equation}\label{tp-truncated-average}
 G_c((\sigma,h),(\sigma',h'))
       =\int_\Omega N_{m+1}(h,h';z,z)
                             \psi(z(\sigma-\sigma'))\,dz.
\end{equation}
The unrestricted Gram matrix is
$G((\sigma,h),(\sigma',h'))=U_m(h,h';t_\sigma,t_{\sigma'})$.
Both are finite sums of positive Gram matrices, with
$0\le G_c\le G$. Finiteness follows also directly from
Theorem~\ref{tp-parseval}. Thus Lemma~\ref{th-cut-contraction}
realizes them as $\langle u_{\sigma,h},u_{\sigma',h'}\rangle$
and $\langle Tu_{\sigma,h},Tu_{\sigma',h'}\rangle$ for a
contraction $T$.

By \eqref{tp-unrestricted} and the inverse-parameter support
proved earlier, these vectors are orthogonal when their
$\rho$ indices differ: a nonzero entry would imply
$\sigma-\sigma'\in R$, contrary to the choice of $A$.
Use the convention that the inner product is antilinear in
its first variable, and define
\[
 v_{\rho,z,h}=\frac1b\sum_{v\in B}\psi(zv)u_{\rho+v,h},\qquad
 v_{z,h}=\frac1a\sum_{\rho\in A}\psi(z\rho)v_{\rho,z,h},
\]
\[
 \widetilde\Theta(z)_{hh'}=
       \frac1a\sum_{\rho\in A}
             \langle v_{\rho,z,h},v_{\rho,z,h'}\rangle.
\]
These expressions are considered for $z\in L_1^\perp$, and
$\widetilde\Theta$ is extended by zero to the other classes
of $K/R$. The values of $\widetilde\Theta$ are $R$-periodic,
since $v\in R$ and $\psi$ is trivial on $R$.

For $z_1\in\Omega$, substitute \eqref{tp-truncated-average}
in $a\langle Tv_{z_1,h},Tv_{z_1,h'}\rangle$. The resulting
factor in the integrand is
\[
 \frac{a}{(ab)^2}
  \left|\sum_{\sigma\in A+B}\psi((z-z_1)\sigma)\right|^2.
\]
The set $A+B$ is a system of representatives of $L_2/L_1$.
Because $z,z_1\in L_1^\perp$, character orthogonality makes
this factor $a$ on $z-z_1\in L_2^\perp$ and zero otherwise.
The whole coset $z_1+L_2^\perp$ lies in $\Omega$ and the
kernel is constant there. Its volume is $1/a$, since
$\operatorname{vol}(L_2)=a$ and the measure is self-dual.
Consequently
\[
 a\langle Tv_{z_1,h},Tv_{z_1,h'}\rangle
                         =N_{m+1}(h,h';z_1,z_1).
\]
Proposition~\ref{th-fourier-average}, with characters
$v\mapsto\psi(-zv)$, now gives
\[
 \widetilde\Theta(z_1)\ge N_{m+1}(z_1,z_1),\qquad
 \sum_{z\in L_1^\perp/R}\widetilde\Theta(z)
   =\frac1{ab}\sum_{\rho\in A}\sum_{v\in B}
                        U_m(t_{\rho+v},t_{\rho+v})
   \le D_{m+1}N_0D_{m+1},
\]
where the final inequality is \eqref{tp-diagonal}.
Replace $\widetilde\Theta$ by $1_{L_c}\widetilde\Theta$.
This keeps positivity and only decreases its total mass.
It still dominates the kernel on $\Omega$, because that
kernel is zero outside $L_c$ by Lemma~\ref{tp-minor-support}.
Every Gram entry between heads with different determinant
absolute values vanishes, since all $|t_\sigma|_K$ are equal.
Thus the required zero entries persist through the averaging.

Finally exhaust $K^\times$ by compact open products of annuli,
stable under all roots of unity. The fixed module $L_c$ works
for every annulus. Proposition~\ref{th-local-limit} gives
one majorant on all units with the same mass and support,
and with the same zero entries. This completes the induction.
If all heads are invertible over $R$, then $c(h)=1$ and
$L_c=R$; the one possible support coset gives the last
inequality directly from the total-mass bound. The finite
target reduction is Proposition~\ref{gu-uniform-germ}.
\end{proof}

The scalar germ, averaged diagonal estimate and all finite
averaging steps of the local induction are therefore proved for
the actual kernel. A local majorant alone does not prove the
four global orbital estimates \eqref{th-four-estimates}; their
uniform arithmetic and automorphic estimates remain distinct
tasks in the complete argument.


